% GENERATED FILE. Edit canonical lessons, then run npm run generate:books.
% canonical-source-hash: fnv1a64:fa51090d116fdac3
% canonical-lessons: KA-C485-mungaru, KA-C485-karmoda, KA-C485-tunturu, KA-C485-jadimale, KA-C485-toyyu

\chapter{The Monsoon Arrives}
\label{ch:ka-the-monsoon-arrives}
\hlchaptermodality{485}

\begin{chapteropening}
\textbf{By the end of this chapter:} \emph{I can say the monsoon, a dark cloud, drizzle, a downpour and to get soaked.}

Complete the last lesson of chapter 485: I can say the monsoon, a dark cloud, drizzle, a downpour and to get soaked.
\end{chapteropening}

\section[\texorpdfstring{\kn{ಮುಂಗಾರು} (muṅgāru)}{ಮುಂಗಾರು (muṅgāru)}]{\kn{ಮುಂಗಾರು} (muṅgāru) — the monsoon; the early rains}

\label{lesson:KA-C485-mungaru}



\begin{quote}
Before the new one: say the Kannada for dampness, then the Kannada for to worsen.
\end{quote}

\subsection*{You'll want to know: \kn{ಮುಂಗಾರು}}
\textbf{\kn{ಮುಂಗಾರು}} — \emph{muṅgāru} — \textquotedblleft{}the monsoon; the early rains\textquotedblright{}.

\begin{quote}
\kn{ಈ} \kn{ವರ್ಷ} \kn{ಮುಂಗಾರು} \kn{ಒಂದು} \kn{ವಾರ} \kn{ತಡವಾಗಿ} \kn{ಬಂದಿತು}.
\end{quote}

\begin{quote}
\emph{ī varṣa muṅgāru ondu vāra taḍavāgi banditu.}
\end{quote}

\begin{quote}
This year the monsoon came a week late.
\end{quote}

The June rains that end the summer heat.

\begin{cousinweb}
\textbf{\kn{ಮುಂಗಾರು}} (\emph{muṅgāru}) is a native Kannada compound of \emph{mun}, \textquotedblleft{}front, first\textquotedblright{} (DEDR 5020), and \emph{kār}, \textquotedblleft{}blackness, the rainy season\textquotedblright{} (DEDR 1278).
\end{cousinweb}

\begin{grammarlens}[title={What you've built}]
One more word for getting around.
\end{grammarlens}

\subsection*{Your turn}
\begin{itemize}
\raggedright
  \item \emph{Say it:} \emph{muṅgāru}
  \item \emph{Say it:} \emph{muṅgāru}, once more
  \item \emph{Say it:} \emph{ī varṣa muṅgāru ondu vāra taḍavāgi banditu.} — This year the monsoon came a week late.
  \item \emph{Recall:} say the Kannada for dampness, then the Kannada for to worsen, then say \emph{muṅgāru} again
\end{itemize}

\begin{tcolorbox}[breakable,colback=teal!4,colframe=teal!35!black,title={Before you move on}]
What does \kn{ಮುಂಗಾರು} mean? (\textquotedblleft{}The monsoon; the early rains\textquotedblright{}.) Say it once more.
\end{tcolorbox}
\section[\texorpdfstring{\kn{ಕಾರ್ಮೋಡ} (kārmōḍa)}{ಕಾರ್ಮೋಡ (kārmōḍa)}]{\kn{ಕಾರ್ಮೋಡ} (kārmōḍa) — a dark cloud; a rain cloud}

\label{lesson:KA-C485-karmoda}



\begin{quote}
Before the new one: say the Kannada for to worsen, then the Kannada for the monsoon.
\end{quote}

\subsection*{You'll want to know: \kn{ಕಾರ್ಮೋಡ}}
\textbf{\kn{ಕಾರ್ಮೋಡ}} — \emph{kārmōḍa} — \textquotedblleft{}a dark cloud; a rain cloud\textquotedblright{}.

\begin{quote}
\kn{ಆಕಾಶ} \kn{ಕಾರ್ಮೋಡಗಳಿಂದ} \kn{ತುಂಬಿತು}.
\end{quote}

\begin{quote}
\emph{ākāśa kārmōḍagaḷinda tumbitu.}
\end{quote}

\begin{quote}
The sky filled with dark clouds.
\end{quote}

Plural: \kn{ಕಾರ್ಮೋಡಗಳು}, \emph{kārmōḍagaḷu}.

\kn{ಕಾರ್}, dark, plus \kn{ಮೋಡ}, a cloud.

\begin{cousinweb}
\textbf{\kn{ಕಾರ್ಮೋಡ}} (\emph{kārmōḍa}) is a native Kannada compound of \emph{kār}, \textquotedblleft{}black, dark\textquotedblright{} (DEDR 1278), and \kn{ಮೋಡ}, \textquotedblleft{}a cloud\textquotedblright{} (DEDR 5033).
\end{cousinweb}

\begin{grammarlens}[title={What you've built}]
One more word for getting around.
\end{grammarlens}

\subsection*{Your turn}
\begin{itemize}
\raggedright
  \item \emph{Say it:} \emph{kārmōḍa}
  \item \emph{Say it:} \emph{kārmōḍa}, once more
  \item \emph{Say it:} \emph{ākāśa kārmōḍagaḷinda tumbitu.} — The sky filled with dark clouds.
  \item \emph{Recall:} say the Kannada for to worsen, then the Kannada for the monsoon, then say \emph{kārmōḍa} again
\end{itemize}

\begin{tcolorbox}[breakable,colback=teal!4,colframe=teal!35!black,title={Before you move on}]
What does \kn{ಕಾರ್ಮೋಡ} mean? (\textquotedblleft{}A dark cloud; a rain cloud\textquotedblright{}.) Say it once more.
\end{tcolorbox}
\section[\texorpdfstring{\kn{ತುಂತುರು} (tunturu)}{ತುಂತುರು (tunturu)}]{\kn{ತುಂತುರು} (tunturu) — drizzle; spray}

\label{lesson:KA-C485-tunturu}



\begin{quote}
Before the new one: say the Kannada for the monsoon, then the Kannada for a dark cloud.
\end{quote}

\subsection*{You'll want to know: \kn{ತುಂತುರು}}
\textbf{\kn{ತುಂತುರು}} — \emph{tunturu} — \textquotedblleft{}drizzle; spray\textquotedblright{}.

\begin{quote}
\kn{ಬೆಳಗ್ಗೆಯಿಂದ} \kn{ತುಂತುರು} \kn{ಮಳೆ} \kn{ಬರುತ್ತಿದೆ}.
\end{quote}

\begin{quote}
\emph{beḷaggeyinda tunturu maḷe baruttide.}
\end{quote}

\begin{quote}
There has been a drizzle since the morning.
\end{quote}

\kn{ತುಂತುರು} \kn{ಮಳೆ}, a light drizzle.

\begin{cousinweb}
\textbf{\kn{ತುಂತುರು}} (\emph{tunturu}) is native Dravidian: DEDR 3398 lists Kannada \emph{tuntuṟ}, \emph{tuntuṟi}, \textquotedblleft{}drizzling, spray, a drop\textquotedblright{}.
\end{cousinweb}

\begin{grammarlens}[title={What you've built}]
One more word for getting around.
\end{grammarlens}

\subsection*{Your turn}
\begin{itemize}
\raggedright
  \item \emph{Say it:} \emph{tunturu}
  \item \emph{Say it:} \emph{tunturu}, once more
  \item \emph{Say it:} \emph{beḷaggeyinda tunturu maḷe baruttide.} — There has been a drizzle since the morning.
  \item \emph{Recall:} say the Kannada for the monsoon, then the Kannada for a dark cloud, then say \emph{tunturu} again
\end{itemize}

\begin{tcolorbox}[breakable,colback=teal!4,colframe=teal!35!black,title={Before you move on}]
What does \kn{ತುಂತುರು} mean? (\textquotedblleft{}Drizzle; spray\textquotedblright{}.) Say it once more.
\end{tcolorbox}
\section[\texorpdfstring{\kn{ಜಡಿಮಳೆ} (jaḍimaḷe)}{ಜಡಿಮಳೆ (jaḍimaḷe)}]{\kn{ಜಡಿಮಳೆ} (jaḍimaḷe) — a downpour; heavy rain}

\label{lesson:KA-C485-jadimale}



\begin{quote}
Before the new one: say the Kannada for a dark cloud, then the Kannada for drizzle.
\end{quote}

\subsection*{You'll want to know: \kn{ಜಡಿಮಳೆ}}
\textbf{\kn{ಜಡಿಮಳೆ}} — \emph{jaḍimaḷe} — \textquotedblleft{}a downpour; heavy rain\textquotedblright{}.

\begin{quote}
\kn{ಜಡಿಮಳೆಯಿಂದ} \kn{ರಸ್ತೆಗಳಲ್ಲಿ} \kn{ನೀರು} \kn{ನಿಂತಿತು}.
\end{quote}

\begin{quote}
\emph{jaḍimaḷeyinda rastegaḷalli nīru nintitu.}
\end{quote}

\begin{quote}
Because of the downpour water stood on the roads.
\end{quote}

\kn{ಜಡಿ}, to pound, plus \kn{ಮಳೆ}, rain.

\begin{cousinweb}
\textbf{\kn{ಜಡಿಮಳೆ}} (\emph{jaḍimaḷe}) is a native Kannada compound of \emph{jaḍi}, \textquotedblleft{}to beat, to pound\textquotedblright{} (DEDR 2300), and \kn{ಮಳೆ}, \textquotedblleft{}rain\textquotedblright{} (DEDR 4753).
\end{cousinweb}

\begin{grammarlens}[title={What you've built}]
One more word for getting around.
\end{grammarlens}

\subsection*{Your turn}
\begin{itemize}
\raggedright
  \item \emph{Say it:} \emph{jaḍimaḷe}
  \item \emph{Say it:} \emph{jaḍimaḷe}, once more
  \item \emph{Say it:} \emph{jaḍimaḷeyinda rastegaḷalli nīru nintitu.} — Because of the downpour water stood on the roads.
  \item \emph{Recall:} say the Kannada for a dark cloud, then the Kannada for drizzle, then say \emph{jaḍimaḷe} again
\end{itemize}

\begin{tcolorbox}[breakable,colback=teal!4,colframe=teal!35!black,title={Before you move on}]
What does \kn{ಜಡಿಮಳೆ} mean? (\textquotedblleft{}A downpour; heavy rain\textquotedblright{}.) Say it once more.
\end{tcolorbox}
\section[\texorpdfstring{\kn{ತೊಯ್ಯು} (toyyu)}{ತೊಯ್ಯು (toyyu)}]{\kn{ತೊಯ್ಯು} (toyyu) — to get soaked; to get drenched}

\label{lesson:KA-C485-toyyu}



\begin{quote}
Before the new one: say the Kannada for drizzle, then the Kannada for a downpour.
\end{quote}

\subsection*{You'll want to know: \kn{ತೊಯ್ಯು}}
\textbf{\kn{ತೊಯ್ಯು}} — \emph{toyyu} — \textquotedblleft{}to get soaked; to get drenched\textquotedblright{}.

\begin{quote}
\kn{ಕೊಡೆ} \kn{ಇಲ್ಲದೆ} \kn{ನಾನು} \kn{ಪೂರ್ತಿ} \kn{ತೊಯ್ದೆ}.
\end{quote}

\begin{quote}
\emph{koḍe illade nānu pūrti toyde.}
\end{quote}

\begin{quote}
Without an umbrella I got soaked through.
\end{quote}

Past (I …): \kn{ತೊಯ್ದೆ}, \emph{toyde}.

\kn{ಒದ್ದೆ} is wet; \kn{ತೊಯ್ಯು} is to become wet through.

\begin{cousinweb}
\textbf{\kn{ತೊಯ್ಯು}} (\emph{toyyu}) is native Dravidian: DEDR 3555 lists Kannada \emph{tuy}, \emph{toy}, \emph{toyyu}, \textquotedblleft{}to be or become wet or moist\textquotedblright{}.
\end{cousinweb}

\begin{grammarlens}[title={What you've built}]
One more word for getting around.
\end{grammarlens}

\subsection*{Your turn}
\begin{itemize}
\raggedright
  \item \emph{Say it:} \emph{toyyu}
  \item \emph{Say it:} \emph{toyyu}, once more
  \item \emph{Say it:} \emph{koḍe illade nānu pūrti toyde.} — Without an umbrella I got soaked through.
  \item \emph{Recall:} say the Kannada for drizzle, then the Kannada for a downpour, then say \emph{toyyu} again
\end{itemize}

\begin{tcolorbox}[breakable,colback=teal!4,colframe=teal!35!black,title={Before you move on}]
What does \kn{ತೊಯ್ಯು} mean? (\textquotedblleft{}To get soaked; to get drenched\textquotedblright{}.) Say it once more.
\end{tcolorbox}
